\documentclass[10pt,a4paper]{amsart}
\usepackage[margin=19mm]{geometry}
\usepackage{amsmath,amssymb,mathtools}
\usepackage[scr=rsfs,cal=euler]{mathalfa}
\usepackage{xfrac}
\usepackage[unicode,hidelinks,bookmarks=false]{hyperref}
\newcommand{\ie}{i.e.\ }
\newcommand{\cf}{cf.\ }
\newcommand{\resp}{respectively\ }
\newcommand{\Subsection}{\subsection}
\providecommand{\nobreakdash}{\nobreak\hbox{-}}
\setlength{\emergencystretch}{2em}
\multlinegap0pt
\usepackage{tikz}
\usetikzlibrary{cd}
\usepackage{mathtools}
\usepackage{amssymb}
\usepackage{enumerate}
\usepackage{xcolor}
\usepackage{float}
\newtheorem{theorem}{Theorem}
\newtheorem{remark}{Remark}
\DeclareMathOperator{\Aut}{Aut}
\newcommand{\calC}{\mathcal{C}}
\newcommand{\calM}{\mathcal{M}}
\newcommand{\mell}{\mathcal{M}_{\text{ell}}}
\newcommand{\mellc}{\overline{\mathcal{M}}_{\text{ell}}}
\newcommand{\nor}{\calC_{E_1 \times_{\mathbb{P}^1} E_2}}
\title{Intersections of the automorphism and the Ekedahl-Oort strata in $M_2$}
\author{Alvaro Gonzalez-Hernandez}
\date{Source: arXiv:2507.07278v2 (CC BY 4.0)}
\begin{document}
\maketitle

\noindent\emph{Source and licence.} Intersections of the automorphism and the Ekedahl-Oort strata in $M_2$. Version arXiv:2507.07278v2 (CC BY 4.0), distributed by arXiv under the Creative Commons Attribution 4.0 International licence, http://creativecommons.org/licenses/by/4.0/, as recorded in the arXiv licence metadata for this exact version and shown on the abstract page https://arxiv.org/abs/2507.07278v2. Full licence notice: \texttt{licenses/arxiv-2507.07278-CC-BY-4.0.txt}.\par\smallskip

\noindent\textit{Editorial scope.} Counted unit: the article's theorem iso\_theorem (source0.tex lines 722--724). The article's complete original proof of it is delivered with it (source0.tex lines 726--883, one \texttt{proof} environment of 158 lines). No mathematics is written, repaired or re-derived: the statement and the proof are the article's own lines, in the article's own notation, and no retained line is substituted or re-flowed.  Retained uncounted as the source-local context the proof needs: lines 590--590.  Nothing was omitted from any retained span, so the package carries no omission record.  No retained line contains a citation, so the wrapper carries no bibliography and no bibliography entry.  The retained lines contain the article's own diagram code, which the wrapper typesets with the standard tikz packages; no image file is delivered and the package declares no asset dependency.  Editorial formatting only, in addition: an AMS \texttt{amsart} preamble that restates the article's own declarations and macros used by the retained lines, namely the environments \texttt{theorem}, \texttt{remark} on separate counters, together with the standard packages those lines need.  The retained environments print in this document as Theorem 1, Remark 1 in order of appearance, whereas the article prints the same items with the numbering of its own full document.  The complete native source of the article as distributed in the arXiv e-print for this exact version is bundled unchanged as \texttt{sources/181331/source0.tex}.
\subsubsection{Parametrisation of the stratum} \label{par_subsection}

\begin{theorem} \label{iso_theorem}
There is an isomorphism $\Phi$ between $\calM$ and $\mell$.
\end{theorem}

\begin{proof}
For any $\calC\in W_{\geq C_2^2}$ and $\sigma$ a non-hyperelliptic involution, there are ten points that are fixed by some non-trivial element of the group $\langle\iota,\sigma\rangle$:
\begin{itemize}
    \item There are six points fixed by $\iota$, namely the Weierstrass points of the curve.  These points are not fixed by $\sigma$ or by $\iota\sigma$, but the action of $\sigma$ permutes them: applying $\sigma$ to a Weierstrass point yields a different Weierstrass point. Therefore, the six Weierstrass points decompose into three orbits under the action of $\sigma$, each consisting of two elements.
    \item There are two points fixed by $\sigma$, which are not fixed by $\iota$ or $\iota\sigma$.
    \item There are two points fixed by $\iota\sigma$, which are not fixed by $\iota$ or $\sigma$.
\end{itemize}


Recall that given a genus two curve, we can construct morphisms $\pi_1:\calC\rightarrow E_1$ and $\pi_2:\calC\rightarrow E_2$, where $E_1$ is the elliptic curve $\calC/\langle\sigma\rangle$
with the choice of the point at infinity given by the image by $\pi_1$ of the fixed points of $\iota\sigma$, and $E_2$ is the elliptic curve $\calC/\langle\iota\sigma\rangle$
with the choice of the point at infinity given by the image by $\pi_2$ of the fixed points of $\sigma$.\\


\underline{\textbf{The morphism $\Phi:\calM\rightarrow\mell$.}}\\

Let $(\calC,\sigma)\in\calM$. As we just mentioned, the action of $\iota$ on $\calC$ fixes the six Weierstrass points, and they form three orbits under the action of $\sigma$. Let us fix an ordering of these three orbits: 
\begin{align*}
\{\{P_{11},P_{12}\},\{P_{21},P_{22}\},\{P_{31},P_{32}\}\},
\end{align*}
where the $P_{ij}$ are the Weierstrass points.\\

Now, let $E_1=\calC/\langle\sigma\rangle$ and $E_2=\calC/\langle\iota\sigma\rangle$ as before. The action of $\iota$ on $\calC$ induces an action of order two on $E_i$, which  is the one that sends each point to its inverse with respect to the group law. The fixed points are the point of infinity and
\begin{align*}
Q_{i,1}&\coloneqq\pi_i(P_{11})=\pi_i(P_{12}), & Q_{i,2}&\coloneqq\pi_i(P_{21})=\pi_i(P_{22}), &
Q_{i,3}&\coloneqq\pi_i(P_{31})=\pi_i(P_{32}).
\end{align*}
Therefore, these points correspond to the non-trivial $2$-torsion points of $E_i$. For $i\in\{1,2\}$, we will then define $\phi_i$ to be the isomorphism:
\begin{align*}
\phi_i\colon\quad(\mathbb{Z}/2\mathbb{Z})^2&\longrightarrow E_i[2](k)\\
(1,0)&\longmapsto Q_{i,1}\\
(0,1)&\longmapsto Q_{i,2}\\
(1,1)&\longmapsto Q_{i,3}.\\
\end{align*}
Suppose that we picked another ordering of the orbits of the Weierstrass points. Let $\tilde{Q}_{i,j}$ the points on $E_i$ associated to this order and $\tilde{\phi_i}$ the corresponding level 2 structure. Then, there exists a $\tau\in S_3$ such that $Q_{i,\tau(j)}=\tilde{Q}_{i,j}$ and, as such, a $\tau\in \Aut((\mathbb{Z}/2\mathbb{Z})^2)$, such that $\tilde{\phi}_i=\phi_i\circ\tau$.\\

Furthermore, if we have two isomorphic pairs $(\calC,\sigma)$ and $(\calC',\sigma')$, there is an isomorphism $\alpha:\calC\rightarrow\calC'$ which induces isomorphisms $\calC/\langle\sigma\rangle\cong\calC'/\langle\alpha\sigma\alpha^{-1}\rangle=\calC'/\langle\sigma'\rangle$ and $\calC/\langle\iota\sigma\rangle\cong\calC'/\langle\alpha\iota\sigma\alpha^{-1}\rangle=\calC'/\langle\iota'\sigma'\rangle$. As $\alpha$ sends the fixed points of $\sigma, \iota\sigma$ and $\iota$ to the fixed points of $\sigma', \iota'\sigma'$ and $\iota'$ respectively, the associated elliptic curves satisfy that $(E_i,\phi_i)\cong(E'_i,\phi_i')$, for $i\in\{1,2\}$.\\

Therefore, the map
\begin{align*}
\Phi\colon\quad\calM&\longrightarrow\mellc\cong (X(2)\times X(2))/\Aut((\mathbb{Z}/2\mathbb{Z})^2)\\
(\calC,\sigma)&\longmapsto((E_1,\phi_1),(E_2,\phi_2))
\end{align*}
is well-defined. The image of $\Phi$ lies in $\mell$ for the following reason:\\

Given an elliptic curve $E_i$, we can always construct a degree two map that ramifies at the four $2$-torsion points. As we can map uniquely any two sets of three points in $\mathbb{P}^1$ using a projective transformation, we deduce that for each $E_i$ there exists a unique map $\psi_i: E_i\rightarrow\mathbb{P}^1$ satisfying
\begin{align*}
\psi_i(\phi_i((1,0)))&=0, & \psi_i(\phi_i((0,1)))&=1, &  \psi_i(\phi_i((1,1)))&=\infty.
\end{align*}
Composing $\psi_i$ with the quotient map $\pi_i:\calC\rightarrow E_i$ gives us maps $\psi_i\circ\pi_i:\calC\rightarrow\mathbb{P}^1$ which satisfy that $\psi_1\circ\pi_1=\psi_2\circ\pi_2$.\\

If $(E_1,\phi_1)\cong(E_2,\phi_2)$, this would imply that there is a $\lambda\in\mathbb{A}^1\setminus\{0,1\}$ and two isomorphisms $\alpha_i:E_i\rightarrow E_\lambda$ such that $\phi_\lambda=\alpha_i\circ\phi_i$. 
One can check that if we define $\psi'$ to be the map $E_\lambda\rightarrow\mathbb{P}^1$ such that
\begin{align*}
\psi'(\phi_\lambda((1,0)))&=0, & \psi'(\phi_\lambda((0,1)))&=1, &  \psi'(\phi_\lambda((1,1)))&=\infty,
\end{align*}
then $\psi'=\psi_1\circ\alpha_1^{-1}=\psi_2\circ\alpha_2^{-1}$, and so,
\begin{align*}
\psi'\circ\alpha_1\circ\pi_1=\psi_1\circ\alpha_1^{-1}\circ\alpha_1\circ\pi_1=\psi_1\circ\pi_1=\psi_2\circ\pi_2=\psi_2\circ\alpha_2^{-1}\circ\alpha_2\circ\pi_2=\psi'\circ\alpha_2\circ\pi_2.\\
\end{align*}
Let $Q$ and $\iota(Q)$ be the two fixed points of $\iota\sigma$ in $\calC$. Then, $\alpha_1\circ\pi_1$ would send both points to the point at infinity of $E_\lambda$, which is a ramification point of $\psi'$. From the description of the fixed points that we gave at the start of the proof, one can check that $\alpha_2\circ\pi_2$ would send $Q$ and $\iota(Q)$ to two different points of $E_\lambda$, none of which is a ramification point of $\psi'$. As $\psi'$ has degree two, we deduce that 
\begin{align*}
\psi'\circ\alpha_1\circ\pi_1(Q)&\neq\psi'\circ\alpha_2\circ\pi_2(Q) & \psi'\circ\alpha_1\circ\pi_1(\iota(Q))&\neq\psi'\circ\alpha_2\circ\pi_2(\iota(Q)) 
\end{align*}
and this would lead to a contradiction, as we had shown that $\psi'\circ\alpha_1\circ\pi_1=\psi'\circ\alpha_2\circ\pi_2$.\\[10pt]

\underline{\textbf{The morphism $\Psi:\mell\rightarrow\calM$.}}\\

Let $((E_1,\phi_1),(E_2,\phi_2))\in X(2)\times X(2)$ such that $(E_1,\phi_1)\not\cong(E_2,\phi_2)$. To construct the inverse of $\Phi$, consider the two morphisms $\psi_1$ and $\psi_2$ that we defined before, which, as we saw, depended on the $\phi_i$.  We define the \textbf{fibre product} of $(E_1,\phi_1)$ and $(E_2,\phi_2)$ along $\mathbb{P}^1$ to be the scheme $E_1\times_{\mathbb{P}^1}E_2$, together with morphisms $p_1$ and $p_2$ to $E_1$ and $E_2$ respectively that satisfies the following universal property: for any scheme $W$
with morphisms $\varphi_1$ and $\varphi_2$, such that $\psi_1\circ\varphi_1=\psi_2\circ\varphi_2$, there exists a unique morphism $\theta:W\rightarrow E_1\times_{\mathbb{P}^1}E_2$ such that $\varphi_1=p_1\circ\theta$ and $\varphi_2=p_2\circ\theta$.

\[\begin{tikzcd}
W \arrow[rd, "\exists\theta" description] \arrow[rrd, "\varphi_1"] \arrow[rdd, "\varphi_2"'] &                                                                &                          \\
                                                                                      & E_1\times_{\mathbb{P}^1}E_2 \arrow[r, "p_1"'] \arrow[d, "p_2"] & E_1 \arrow[d, "\psi_1"'] \\
                                                                                      & E_2 \arrow[r, "\psi_2"]                                        & \mathbb{P}^1            
\end{tikzcd}\vspace{10pt}\]

We will prove that the normalisation $\nor$ of $E_1\times_{\mathbb{P}^1}E_2$ is a genus two curve. To do this, we will show that for any $E_1$ and $E_2$, there is a genus two curve $\calC_{\lambda_1,\lambda_2}$ and morphisms $\varphi_1$ and $\varphi_2$ making the diagram commutative and use the information that we know about the ramification of the morphisms to show that this curve is isomorphic to $\nor$.\\

Let $\mu_i\in\mathbb{P}^1$ be the image under $\psi_i$ of the point at infinity of $E_i$, and let $\lambda_i=\tfrac{1}{1-\mu_i}$. As we saw, we can construct an isomorphism $\alpha_i:E_i\rightarrow E_{\lambda_i}$ where
\begin{align*}
 E_{\lambda_i}\colon\quad y^2=x(x-1)(x-\lambda_i)
\end{align*}
such that $(E_i,\phi_i)\cong(E_{\lambda_i},\phi_{\lambda_i})$ and $\psi_i=\psi_i'\circ\alpha_i$ is defined by
\begin{align*}
\psi_i'\colon\quad E_{\lambda_i}&\longrightarrow\mathbb{P}^1\\
(x,y)&\longmapsto\frac{(1-\lambda_i)x}{x-\lambda_i}\\
\end{align*}
Recall that $(E_1,\phi_1)\not\cong(E_2,\phi_2)$ implies that $\lambda_1\neq\lambda_2$. Now, consider the curve
\begin{align*}
\calC_{\lambda_1,\lambda_2}\colon\quad y^2=(x^2-1)\left(x^2-\frac{\lambda_1}{\lambda_2}\right)\left(x^2-\frac{\lambda_1(\lambda_2-1)}{\lambda_2(\lambda_1-1)}\right).\\
\end{align*}
The curve $\calC_{\lambda_1,\lambda_2}$ has genus two whenever its discriminant is non-zero. This occurs precisely for any two values $\lambda_1,\lambda_2\in \mathbb{A}^1\setminus\{0,1\}$, such that $\lambda_1\neq\lambda_2$. We can construct two maps, $\rho_1$ and $\rho_2$
\begin{align*}
\rho_1\colon\quad \calC_{\lambda_1,\lambda_2} &\longrightarrow E_{\lambda_1} \\
(x,y) &\longmapsto\left(\nu_1^2\left(x^2-\frac{\lambda_1(\lambda_2-1)}{\lambda_2(\lambda_1-1)}\right),\nu_1^3 y\right),\\
\rho_2\colon\quad \calC_{\lambda_1,\lambda_2} &\longrightarrow E_{\lambda_2} \\
(x,y) &\longmapsto\left(-\nu_2^2\left(\frac{1}{x^2}-\frac{\lambda_2(\lambda_1-1)}{\lambda_1(\lambda_2-1)}\right),\nu_2^3\left( \frac{\lambda_2(\lambda_1-1)^{\frac{1}{2}}y}{\lambda_1(\lambda_2-1)^{\frac{1}{2}}x^3}\right)\right),
\end{align*}

where
\begin{align*}
\nu_1&=\left(\frac{\lambda_2(\lambda_1-1)}{\lambda_1-\lambda_2}\right)^\frac{1}{2}, & \nu_2&=\left(\frac{\lambda_1(\lambda_2-1)}{\lambda_1-\lambda_2}\right)^\frac{1}{2}.\\
\end{align*}
Then, $\psi_1'\circ\rho_1=\psi_2'\circ\rho_2$, as in both cases, this is the map
\begin{align*}
\psi_i'\circ\rho_i\colon\quad \calC_{\lambda_1,\lambda_2}&\longrightarrow \mathbb{P}^1\\
(x,y)&\longmapsto \frac{\lambda_2(1-\lambda_1)x^2-\lambda_1(1-\lambda_2)}{\lambda_2x^2-\lambda_1}.\\
\end{align*}
As a consequence, we can define  $\varphi_1=\alpha_1^{-1}\circ\rho_1$ and $\varphi_2=\alpha_2^{-1}\circ\rho_2$, and deduce that there exists a map $\theta:\calC_{\lambda_1,\lambda_2}\rightarrow E_1\times_{\mathbb{P}^1}E_2$ making the following diagram commutative:
\[
\begin{tikzcd}
\calC_{\lambda_1,\lambda_2} \arrow[rd, "\theta" description] \arrow[rrd, "\varphi_1"] \arrow[rdd, "\varphi_2"'] \arrow[rr, "\rho_1"] \arrow[dd, "\rho_2"'] &                                                                & E_{\lambda_1} \arrow[dd, "\psi_1'" description, bend left] \\
                                                                                                                                        & E_1\times_{\mathbb{P}^1}E_2 \arrow[r, "p_1"'] \arrow[d, "p_2"] & E_1 \arrow[d, "\psi_1"'] \arrow[u, "\alpha_1"]               \\
E_{\lambda_2} \arrow[rr, "\psi_2'" description, bend right]                                                                             & E_2 \arrow[r, "\psi_2"] \arrow[l, "\alpha_2"']                   & \mathbb{P}^1                                              
\end{tikzcd}\vspace{10pt}
\]
It is not difficult to see that the preimages of the three points that are simultaneously in the branch loci of $\psi_1$ and $\psi_2$ give rise to singularities of $E_1\times_{\mathbb{P}^1}E_2$, so this is not a smooth curve. As $p_1$ is a dominant morphism to $E_1$ of degree two, $E_1\times_{\mathbb{P}^1}E_2$ has at most two irreducible components.\\

If it had two irreducible components, the ramification locus of $p_1$ should correspond to the points where the two components meet and, therefore, all of these points should be singular. However, we can check that $p_1$ has two ramification points which are smooth, namely, the two preimages of the ramification point of $\psi_1$ that is not a ramification point of $\psi_2$ (earlier in the proof we denoted this point as $\mu_1$). Therefore, we deduce that $E_1\times_{\mathbb{P}^1}E_2$ is irreducible. Although we will not prove it, one can further check that $E_1\times_{\mathbb{P}^1}E_2$ has three nodal singularities, and its arithmetic genus is five.\\

As both maps $\varphi_1$ and $p_1$ have degree two, we deduce that the map $\theta$ is generically $1$-to-$1$, and as $\calC_{\lambda_1,\lambda_2}$ is smooth, we deduce that $\calC_{\lambda_1,\lambda_2}$ is isomorphic to the normalisation of $E_1\times_{\mathbb{P}^1}E_2$. There is a clear non-hyperelliptic involution
in $\calC_{\lambda_1,\lambda_2}$ sending $x\mapsto -x$, and this gives rise to an involution $\tilde{\sigma}$ in $\nor$.
Therefore, we have a morphism
\begin{align*}
(X(2)\times X(2))\setminus\Delta&\longrightarrow\calM\\
((E_1,\phi_1),(E_2,\phi_2))&\longmapsto(\nor,\tilde{\sigma})\\
\end{align*}
This extends to a morphism $\Psi:\mell\rightarrow\calM$.\\

To see this, assume that we have $((E'_1, \phi'_1), (E'_2, \phi'_2))$ such that $E_1 \cong E'_1$, $E_2 \cong E'_2$, and there is an automorphism $\tau \in \Aut((\mathbb{Z}/2\mathbb{Z})^2)$ with $\tau(\phi_i) = \phi'_i$ for $i = 1, 2$.  
Then, from the description of $X(2)$ that we gave earlier, we see that $\mathcal{C}_{E'_1 \times_{\mathbb{P}^1} E'_2} \cong \mathcal{C}_{f_\tau(\lambda_1), f_\tau(\lambda_2)}$.  
By computing the Igusa invariants with Magma, we can check that 
\[
J_{2n}(\mathcal{C}_{\lambda_1, \lambda_2}) = J_{2n}(\mathcal{C}_{f_\tau(\lambda_1), f_\tau(\lambda_2)})
\]
for all $1 \leq n \leq 5$ and all $\tau \in S_3$. Therefore, $\mathcal{C}_{\lambda_1, \lambda_2} \cong \mathcal{C}_{f_\tau(\lambda_1), f_\tau(\lambda_2)}$.\\

Hence, $\Psi: \mathcal{M}_{\mathrm{ell}} \to \mathcal{M}$ is a well-defined morphism.

\begin{remark}
While it is always true that $\mathcal{C}_{\lambda_1, \lambda_2} \cong \mathcal{C}_{f_\tau(\lambda_1), f_\tau(\lambda_2)}$, it is certainly not true in general that $\mathcal{C}_{\lambda_1, \lambda_2} \cong \mathcal{C}_{\lambda_1, f_\tau(\lambda_2)}$.  
In fact, one can check that $\mathcal{C} \in W_{\geq D_4}$ if and only if there exists $\lambda \in \mathbb{A}^1 \setminus \{0,1\}$ and an involution $\tau \in S_3$ such that $\mathcal{C} \cong \mathcal{C}_{\lambda, f_\tau(\lambda)}$.  
Likewise, $\mathcal{C} \in Z$, where $Z$ is as defined in subsection~\ref{par_subsection}, if and only if there exists $\lambda$ and an element $\tau \in S_3$ of order three such that $\mathcal{C} \cong \mathcal{C}_{\lambda, f_\tau(\lambda)}$.\\
\end{remark}
\newpage

\underline{$\Phi=\Psi^{-1}$}\\

Finally, to prove that we have an isomorphism $\Phi:\calM\rightarrow\mell$, we need to check that $\Psi\circ\Phi=\mathrm{id}_{\calM}$. This is easy to see from the following commutative diagram:
\[\begin{tikzcd}
\calC \arrow[rrd, "\pi_1"] \arrow[rdd, "\pi_2"'] &                                                                                                                   &                                                 \\
                                                 & \quad\calC/\langle\sigma\rangle\!\times_{\mathbb{P}^1}\!\calC/\langle\iota\sigma\rangle \arrow[r, "p_1"'] \arrow[d, "p_2"] & \calC/\langle\sigma\rangle \arrow[d, "\psi_1"'] \\
                                                 &\calC/\langle\iota\sigma\rangle \arrow[r, "\psi_2"]                                                               & \mathbb{P}^1                                   
\end{tikzcd}\]

The universal property of the fibre product implies that we can construct a map $\theta:\calC\rightarrow\calC/\langle\sigma\rangle\times_{\mathbb{P}^1}\calC/\langle\iota\sigma\rangle$. This map $\theta$ induces an isomorphism between $\calC$ and the normalisation of $\calC/\langle\sigma\rangle\times_{\mathbb{P}^1}\calC/\langle\iota\sigma\rangle$. By carefully tracking the involution $\sigma$ through the commutative diagrams, one can check that $\Psi\circ\Phi=\mathrm{id}_{\calM}$.
\end{proof}

\end{document}
